The accretion of matter onto compact objects, such as neutron stars and black
holes, is one of the most efficient ways to produce radiant energy in
astrophysical systems. Consider an element of gas of mass $\Delta m$ in a
stationary and geometrically thin disc of matter with a maximum radius of
$R_{\mathrm{max}}$ and minimum stable orbital radius of $R_{\mathrm{min}}$
(with $R_{\mathrm{min}}/R_{\mathrm{max}} \ll 1$), in rotation around a compact
object of mass $M$ and radius $R$.

\begin{description}
\item[(a)] Assuming that an element of gas in the disc follows an
  approximately Keplerian circular orbit, find an expression for the total
  mechanical energy per unit mass $\frac{\Delta E}{\Delta m}$ released by this
  gas from when it starts orbiting at a radius $R_{\mathrm{max}}$ until it
  reaches an orbital radius of $r \ll R_{\mathrm{max}}$. This process occurs
  very slowly, transforming kinetic energy into internal energy of the gas
  disc through viscous dissipation. \textbf{Note:} Ignore the gravitational
  interaction between particles within the accretion disc and give your final
  answer in terms of $G$, $M$, and $r$. \hfill[6\,pt]
\item[(b)] Considering that the disc receives mass at an average rate of
  $\dot{M}$, and assuming that all the mechanical energy lost is ultimately
  converted into radiation, find an expression for the total luminosity of the
  disc. \hfill[5\,pt]
\item[(c)] Consider now the ring composed of all mass elements with radii
  between $r$ and $r + \Delta r$. In this scenario, find an expression for the
  luminosity generated by the disc over its small width $\Delta r$ at this
  radius; that is, find the expression for
  $\frac{\Delta E}{\Delta t\,\Delta r}$. \hfill[8\,pt]
\item[(d)] Assuming that the gravitational energy released in this ring is
  locally emitted by the surface of the ring in the form of black-body
  radiation, find an expression for the surface temperature $T$ of the ring.
  \hfill[10\,pt]
\item[(e)] Consider that the central object is a stellar black hole with a
  mass of $3M_\odot$ and a rate of accretion of
  $\dot{M} = 10^{-9}\,M_\odot/\mathrm{year}$. Consider also that
  $R_{\mathrm{min}} = 3R_{\mathrm{sch}}$, where $R_{\mathrm{sch}}$ is the
  Schwarzschild radius of the black hole. Determine the luminosity of the disc
  and the peak wavelength of emission of its innermost part. Ignore
  gravitational redshift effects and assume that the emission from the
  innermost part of the ring dominates the total emission. \hfill[3\,pt]
\item[(f)] Now, considering another accretion system with
  $\dot{M} = 1\,M_\odot/\mathrm{year}$ and a peak emission wavelength of
  $\lambda = 6 \times 10^{-8}$\,m, estimate the mass of this black hole.
  \hfill[3\,pt]
\end{description}
